\section{Discussion}\label{sec:discussion}

\subsection{What the theorem says and does not say}\label{sec:discussion-scope}

Theorem~\ref{thm:main} is a statement about the limit $q\to\infty$. It gives no bound for
any particular $q$, and none can be extracted from the proof at a size of practical
interest.

\emph{The thresholds in the proof are astronomically large.} In the comparison between the
market and the independent-team model, the error term is negligible only when $\ln q$
exceeds an explicit multiple of $1/\tau$, where $\tau$ is one of the small parameters of the
error budget; with the parameters certified for $\beta=2.30$ this means $q$ above about
$10^{52{,}900}$, and for $\beta=2.302$ the parameter $\tau$ is twenty times smaller
(Remark~\ref{rem:transfer-size}). The local picture behind the fractional upper bound also sets in late: a
team of rescaled cost at most $K$ shares a pair with about $\tfrac{9}{10}K^2q^{-1/3}$ other
such teams, and for $K=6$ this number stays above $1$ until $q$ is about $3\cdot10^4$. Below
such sizes the near-tight hypergraph of the market is far from a tree.

\emph{The constants of the proof are not visible at sizes that can be simulated.} In the
market, the measured value of $\PhiLP/q^{1/3}$ is between $2.53$ and $2.56$ for $q\le48$, and at $q=2000$
the fractional optimum restricted to teams of rescaled cost at most $6$, which is an upper
bound on $\PhiLP$, is still $2.31\,q^{1/3}$ (Table~\ref{tab:discussion-market}; raising the
cut-off to $7$ changed it by about $10^{-4}$). That is far from the limiting bound
$\alphalp=2.2519$. So the two inequalities
that the proof compares, $\Phiq\ge\betaint q^{1/3}$ and $\E\PhiLP\le(\alphalp+o(1))q^{1/3}$,
do not describe the market at any size we could compute.

\emph{The conclusion, by contrast, is visible early.} The measured probability of a nonempty
core is $0.37$ at $q=16$, $0.12$ at $q=24$, $0.04$ at $q=32$ and $0.002$ at $q=48$
(Table~\ref{tab:discussion-core}), and the measured mean gap grows over the same range. The
theorem therefore agrees with the data in its conclusion, while its proof describes a regime
that begins many orders of magnitude later. Nothing proved here explains the behaviour at
small $q$.

\emph{Further limits.}
\begin{itemize}
\item The constant $0.050$ in part (i) is the difference of two certified bounds and is not
claimed to be sharp. In the independent-team model the measured gap is about $0.09$ to
$0.115$ times $q^{1/3}$ for $q\le40$ and rising.
\item The proof gives only a polynomial rate in part (iii): by Chebyshev's inequality, part
(i) and Theorem~\ref{thm:variance}, $\P(\Dq=0)\le\Var\Dq/(\E\Dq)^2=O(q^{-5/6+\eps})$.
\item The theorem is for independent uniform pair values, three sectors of equal size, and
teams that need all three pairs.
\item The proof imports Theorem~2(b) of \cite{FPT20}, and two of its steps are
computer-assisted (Section~\ref{sec:certificates}).
\item Every part of the proof was written and checked by AI agents. In the working notes
from which this text was consolidated, each component was attacked by at least two further
agents of the same model family and read at least once by a model of a different family, and
no error affecting a theorem was found. No person has checked any part.
\end{itemize}

\subsection{Numerical evidence}\label{sec:discussion-numerics}

\emph{Everything in this subsection is floating-point simulation. None of it is used in any
proof, and none of it has been checked by a person.} Linear programmes were solved in double
precision; integer optima were computed by branch and bound and are reported only where the
solver closed the instance. Standard errors are given where the data files record them.

\begin{table}[ht]
\caption{Floating point. The market: share of instances with a nonempty core, with $95\%$
Clopper--Pearson intervals. The core was declared nonempty when the vertex returned by the
linear-programming solver was integral; for $q\le24$ a fractional vertex was followed by an
exact integer solve and a comparison of the two optima. Source:
\protect\path{gate1_ext/out/ext_table_and_fits.txt}.}\label{tab:discussion-core}
\begin{tabular}{@{}rrrcc@{}}
\toprule
$q$ & instances & nonempty & share & $95\%$ interval \\
\midrule
2 & 2000 & 2000 & 1.000 & $[0.998,1.000]$ \\
3 & 2000 & 1987 & 0.994 & $[0.989,0.997]$ \\
4 & 2000 & 1939 & 0.970 & $[0.961,0.977]$ \\
5 & 2000 & 1873 & 0.936 & $[0.925,0.947]$ \\
6 & 2000 & 1784 & 0.892 & $[0.878,0.905]$ \\
8 & 2000 & 1568 & 0.784 & $[0.765,0.802]$ \\
10 & 2000 & 1338 & 0.669 & $[0.648,0.690]$ \\
12 & 1000 & 582 & 0.582 & $[0.551,0.613]$ \\
16 & 1000 & 366 & 0.366 & $[0.336,0.397]$ \\
20 & 1000 & 211 & 0.211 & $[0.186,0.238]$ \\
24 & 500 & 61 & 0.122 & $[0.095,0.154]$ \\
28 & 500 & 40 & 0.080 & $[0.058,0.107]$ \\
32 & 500 & 19 & 0.038 & $[0.023,0.059]$ \\
40 & 2000 & 16 & 0.008 & $[0.005,0.013]$ \\
48 & 1000 & 2 & 0.002 & $[0.000,0.007]$ \\
\bottomrule
\end{tabular}
\end{table}

An independent sample gave shares $0.353$, $0.134$, $0.037$ and $0$ of $211$ at
$q=16,24,32,48$. Among the curves fitted to Table~\ref{tab:discussion-core}, the stretched
exponential $\exp(-0.0115\,(q-2)^{1.68})$ fits best without an extra parameter (deviance
$20.3$ on $13$ degrees of freedom); a pure exponential in $q$ fits badly (deviance $1505$ on
$14$). The fits are descriptive and cover $q\le48$ only.

\begin{table}[ht]
\caption{Floating point. The market: integer and fractional deficits divided by $q^{1/3}$,
and the mean gap. Upper block: both optima on the same instances
(\protect\path{lemmaC/round3/scout/REPORT.md}). Lower block: fractional optimum only,
restricted to teams of rescaled cost at most $6$ or $7$, which is an upper bound on $\PhiLP$
(\protect\path{lemmaC/round3/redteam2/NOTES.md}); the row $q=160$ is from an earlier run,
recorded as $2.46$ and as $2.47$.}\label{tab:discussion-market}
\begin{tabular}{@{}rrcccc@{}}
\toprule
$q$ & instances & $\Phiq/q^{1/3}$ & $\PhiLP/q^{1/3}$ & $\E\Dq$ & $\P(\Dq>0)$ \\
\midrule
16 & 1000 & 2.572 & 2.556 & 0.041 & 0.647 \\
24 & 1000 & 2.586 & 2.561 & 0.072 & 0.866 \\
32 & 400 & 2.598 & 2.565 & 0.104 & 0.963 \\
48 & 211 & 2.579 & 2.535 & 0.161 & 1.000 \\
\midrule
160 & 8 & & 2.46--2.47 & & \\
500 & 1 & & 2.39 & & \\
1000 & 4 & & $2.352\pm0.006$ & & \\
2000 & 3 & & $2.313\pm0.011$ & & \\
\bottomrule
\end{tabular}
\end{table}

In the market the fractional constant keeps falling with $q$. Its excess over $2.2$ is
between about $1$ and $1.5$ times $q^{-1/3}$ from $q=16$ to $q=2000$, which is compatible
with a limit near the tree values and a correction of order $q^{-1/3}$, as the slow
disappearance of pair-sharing suggests. The integer
optimum of the market was not computed beyond $q=48$. In the instances with $\Dq>0$ the
optimal fractional solution found by the solver has non-integral weights at a share $0.71$,
$0.77$, $0.83$, $0.87$ of the agents for $q=16,24,32,48$; so the fractional optimum is not a
small perturbation of an assignment. The standard deviation of $\Dq$ is
about $0.05$ for $q$ between $16$ and $48$, in line with Theorem~\ref{thm:variance}.

\begin{table}[ht]
\caption{Floating point. The independent-team model, teams of rescaled cost at most $6$:
fractional deficit divided by $q^{1/3}$ (mean and standard error), mean gap divided by $q^{1/3}$ on the instances
whose integer optimum was found (all were closed to optimality), their sum as an estimate of
the integer constant, and the share of instances with an integral fractional optimum. Source:
\protect\path{lemmaC/round3/pscout/results.txt}.}\label{tab:discussion-iid}
\small
\begin{tabular}{@{}rrcrccc@{}}
\toprule
 & \multicolumn{2}{c}{fractional} & \multicolumn{2}{c}{gap} & integer & share \\
$q$ & instances & constant & instances & mean & constant & integral \\
\midrule
16 & 3000 & $2.1654\pm0.0025$ & 600 & $0.0878\pm0.0016$ & 2.253 & 0.027 \\
24 & 163 & $2.1757\pm0.0086$ & 163 & $0.1011\pm0.0024$ & 2.277 & 0.006 \\
32 & 2000 & $2.1819\pm0.0020$ & 74 & $0.1029\pm0.0028$ & 2.285 & 0.000 \\
40 & 19 & $2.1960\pm0.0189$ & 19 & $0.1149\pm0.0045$ & 2.311 & 0.000 \\
64 & 1000 & $2.1891\pm0.0022$ & & & & \\
128 & 400 & $2.1933\pm0.0024$ & & & & \\
256 & 200 & $2.1977\pm0.0023$ & & & & \\
512 & 20 & $2.1996\pm0.0044$ & & & & \\
\bottomrule
\end{tabular}
\end{table}

In the independent-team model there is no pair-sharing and the constants are close to their
limits already at small $q$. The fractional constant rises from $2.165$ to $2.200$ and fits
give a limit between $2.20$ and $2.22$; a second code found $2.184\pm0.012$ at $q=1000$,
$2.202$ at $q=2000$ and $2.196$ at $q=4000$ (two and one instances). All of these are below
$\alphalp=2.2519$ and below the floating-point cost $2.2495$ of the matching defined by the
fixed point. This is consistent with $\alphalp$ as an upper bound for the limiting
fractional constant, which Theorems~\ref{thm:lp} and~\ref{thm:existence} prove for the
market; the corresponding statement for the independent-team model is not proved in this
paper. The
integer constant rises from $2.253$ to $2.311$ over $q=16,\dots,40$, and fits give a limit of
$2.32$ or more; a second code found $2.32\pm0.03$ at $q=32$. This is consistent with
$\betaint=2.302$ as a lower bound for the limit. The finite-size values at $q\le32$ lie below
$2.302$; the theorems concern the limit only. The comparison was run as a test that could
have refuted the proof's architecture, and it did not.

\subsection{The matching on the tree is not a function of the tree}
\label{sec:discussion-endogeny}

Theorems~\ref{thm:integer} and~\ref{thm:lp} together have a consequence for the tree itself,
which is proved as Corollary~\ref{cor:proofmain-endogeny}. In words:

\emph{No invariant perfect matching on $\tree$ of cost below $\betaint$ is a measurable
function of the marked tree.} In particular the matching built from a fixed point of $\RDE$
(Theorem~\ref{thm:process}), whose cost is at most $\alphalp<\betaint$, uses randomness
beyond the tree and its costs. In the usual terminology, the recursive tree process attached
to the recursion \eqref{eq:intro-rde} is not endogenous.

The reason is the following. The proof of Theorem~\ref{thm:lp} uses fractionality at one
place only: the weight given to a team is the conditional probability, given a ball of
finite radius, that the team belongs to the matching $m$. If $m$ were a measurable function
of the marked tree, these conditional probabilities would converge to the indicator of $m$
itself as the radius grows, and a thresholded version of the construction (keep a team when
its weight exceeds $1/2$ at all three of its agents, then patch) would give \emph{integral}
assignments of cost $(c(m)+o(1))q^{1/3}$ in the finite market. By
Theorem~\ref{thm:integer} this forces $c(m)\ge\betaint$.

This matches a simulation (floating point): two copies of the process run on the same tree
from independent initial conditions do not merge. Their root variables settle at correlation
about $0.51$, and they choose the same team at the root with probability about $0.41$. It
also explains why the gap is invisible locally: on the tree the fractional and the integer
problem have the same value (Theorem~\ref{thm:decomp}), but no cheap invariant matching of
the tree is a measurable function of the tree and its costs
(Corollary~\ref{cor:proofmain-endogeny}).

\subsection{Open problems}\label{sec:discussion-open}

\emph{The constants.} Do $\Phiq/q^{1/3}$ and $\PhiLP/q^{1/3}$ converge, and to what? For the
integer problem we prove $\liminf\ge2.302$ in probability; the star-conditioned first moment
cannot give more than about $2.3025$, and the simulations suggest a limit of $2.32$ or more.
For the fractional problem we prove $\limsup\E\PhiLP/q^{1/3}\le2.2519$, the fixed point costs
$2.2495$ numerically, and the simulations of the independent-team model suggest a limit near
$2.20$. If that is right, the matching built from the recursion is not the cheapest invariant
matching on $\tree$. Is the fractional limit equal to the infimum of $c(m)$ over invariant
perfect matchings on $\tree$, and what is that infimum? The plain first moment gives the lower bound $3e^{-1/3}=2.1496$ for the integer
problem only; no lower bound for the fractional constant is proved in this paper.

\emph{The rate.} The proof gives $\P(\text{core nonempty})=O(q^{-5/6+\eps})$. The data for
$q\le48$ decay much faster, like a stretched or plain exponential in $q$. What is the true
rate? A Chebyshev argument cannot give more than a polynomial rate.

\emph{Other sector graphs and more sectors.} Here every pair of sectors interacts. If one of
the three types of pair is removed, the sector graph is a path, and for sector graphs that
are trees the core is expected to be nonempty always \todo{cite the deterministic result for
tree sector graphs}. With four or more sectors and pair-additive values on a cyclic sector
graph, nothing is known; the scaling of the optimum for the four-sector model with values on
all pairs is itself not treated in \cite{FPT20}. Unequal sector sizes are a further case:
for independent values on triples, the figures reported in \cite{Lucas95} suggest that the
core is nonempty with probability tending to $1$ when one side stays small.

\emph{Other value laws.} The proof uses the uniform law through the density of a pair
deficit near zero and through the imported bound of \cite{FPT20}. Pair values with a bounded
density that is positive at the top of the support should behave in the same way. Unbounded
values, for which the two-sided law of one price fails \cite{HassidimRomm15}, and values
that are independent over triples \cite{Lucas95,MMR05}, are different problems.

\emph{Removing the computer assistance.} The main theorem needs only some $\beta>\alpha$
with $\Phiq\ge\beta q^{1/3}$ with high probability and $9\,\E X\le\alpha$. The margin is
small: the plain first moment gives $2.1496$, the star-conditioned one $2.3022$, and the
bounds $c_n$ on the tree cost fall below $2.302$ only from the fourth iterate on
($c_4=2.3017$, $c_5=2.2679$ in floating point; the certified bounds are below $2.302$ from
stage $5$ on). A proof without computation would need an analytic bound on a five-fold
iterated integral, or a different argument for either constant. Short of that, the two
certificates should be re-implemented in a second interval library, and the scripts read by
a person.
